% ── §3.1: Περιγραφή των Datasets ──────────────────────────────

\section{Περιγραφή των Datasets}
\label{sec:datasets}

Η πειραματική αξιολόγηση της παρούσας εργασίας στηρίζεται σε δύο
συμπληρωματικά σύνολα δεδομένων: το ευρέως αναφερόμενο UCI HAR
ως σημείο αναφοράς (baseline) σε ελεγχόμενο περιβάλλον, και το
AAL Smartphone Dataset ως κύριο πειραματικό σύνολο σε συνθήκες
πραγματικής χρήσης.

% ─── 3.1.1 UCI HAR ──────────────────────────────────────────
\subsection{UCI Human Activity Recognition Dataset}
\label{subsec:uci_har}

Το \textbf{UCI HAR Dataset} \cite{Anguita2013} αποτελεί το πλέον
ευρέως χρησιμοποιούμενο σύνολο αναφοράς στη βιβλιογραφία HAR.
Συλλέχθηκε από 30 εθελοντές (ηλικίας 19–48 ετών) σε ελεγχόμενο
εργαστηριακό περιβάλλον, φορώντας Samsung Galaxy S II στη μέση.
Ο αισθητήρας δειγματοληπτεί στα 50\,Hz και οι 30 συμμετέχοντες
κατανέμονται σε σύνολο εκπαίδευσης (70\%, 21 άτομα, 7352 παράθυρα)
και σύνολο ελέγχου (30\%, 9 άτομα, 2947 παράθυρα). Τα βασικά
χαρακτηριστικά συνοψίζονται στον Πίνακα~\ref{tab:uci_har}.

\begin{table}[H]
    \centering
    \caption{Βασικά χαρακτηριστικά UCI HAR Dataset.}
    \label{tab:uci_har}
    \begin{tabular}{ll}
        \toprule
        \textbf{Παράμετρος} & \textbf{Τιμή} \\
        \midrule
        Συμμετέχοντες          & 30 (21 train / 9 test) \\
        Ηλικιακό εύρος         & 19–48 ετών \\
        Περιβάλλον             & Ελεγχόμενο (εργαστήριο) \\
        Αισθητήρες             & Επιταχυνσιόμετρο + Γυροσκόπιο \\
        Συχνότητα δειγματοληψίας & 50\,Hz \\
        Μήκος παραθύρου        & 128 δείγματα (2{,}56\,s) \\
        Επικάλυψη              & 50\% \\
        Κανάλια εισόδου        & 9 \\
        Παράθυρα εκπαίδευσης   & 7352 \\
        Παράθυρα ελέγχου       & 2947 \\
        Αριθμός κλάσεων        & 6 \\
        \bottomrule
    \end{tabular}
\end{table}

Κάθε παράθυρο αντιστοιχεί σε ένα διάνυσμα
$\mathbf{X}^{(k)} \in \mathbb{R}^{9 \times 128}$, το οποίο
αποτελείται από: τριαξονική ολική επιτάχυνση ($a_x, a_y, a_z$),
τριαξονική γωνιακή ταχύτητα ($g_x, g_y, g_z$) και τριαξονική
επιτάχυνση σώματος (body acceleration, $b_x, b_y, b_z$) μετά
από αφαίρεση της βαρύτητας με low-pass φίλτρο Butterworth.
Τα δεδομένα παρέχονται ήδη κανονικοποιημένα στο διάστημα
$[-1, 1]$ \cite{Anguita2013}.

Ο διαχωρισμός ανά χρήστη (subject-based split) — εκπαίδευση με
21 subjects (7352 παράθυρα) και έλεγχος με 9 subjects (2947 παράθυρα) —
διασφαλίζει ότι η αξιολόγηση γίνεται σε δεδομένα από
\textit{άγνωστους} χρήστες, αποτυπώνοντας τη γενικευσιμότητα του
μοντέλου. Στο FL σύστημά μας, κάθε subject του training set
αντιστοιχίζεται σε έναν client. Το UCI HAR χρησιμοποιείται ως
\textit{baseline}: δεδομένου ότι η συλλογή έγινε σε ελεγχόμενες
συνθήκες (κάθε εθελοντής εκτέλεσε όλες τις δραστηριότητες σε
καθορισμένη ακολουθία), αναμένουμε σχεδόν IID κατανομή labels
μεταξύ subjects — υπόθεση που επιβεβαιώνεται ποσοτικά στο §3.2
και θέτει το κατώτατο όριο δυσκολίας για FL.

% ─── 3.1.2 AAL Smartphone Dataset ───────────────────────────────
\subsection{AAL Smartphone Dataset}
\label{subsec:aal_dataset}

Το \textbf{AAL Smartphone Dataset} --- το οποίο δημιουργήθηκε από
τους Davis \& Owusu \cite{DavisOwusu2016} και διατίθεται δημόσια μέσω του UCI Machine Learning Repository
\cite{AAL_UCI2016} --- επεκτείνει τα δεδομένα του UCI HAR σε συνθήκες
πραγματικής Υποβοηθούμενης Διαβίωσης. Συλλέχθηκε από 30 συμμετέχοντες
ευρύτερου ηλικιακού φάσματος (22–79 ετών) — συμπεριλαμβανομένων
ηλικιωμένων — με smartphone τοποθετημένο στη μέση, στα 50\,Hz.
Το σύνολο διαχωρίζεται subject-independently σε \textbf{22 χρήστες εκπαίδευσης}
(οι οποίοι λειτουργούν ως FL clients) και 8 χρήστες ελέγχου· καθώς το dataset
δεν παρέχει ρητά αναγνωριστικά συμμετεχόντων (subject IDs), αυτά ανακατασκευάζονται
από τη δομή των block της ακολουθίας δραστηριοτήτων (κάθε χρήστης εκτελεί διαδοχικά
τις 6 δραστηριότητες).
Τα βασικά χαρακτηριστικά του dataset παρατίθενται στον Πίνακα~\ref{tab:aal_dataset}.

\begin{table}[H]
    \centering
    \caption{Βασικά χαρακτηριστικά AAL Dataset.}
    \label{tab:aal_dataset}
    \begin{tabular}{ll}
        \toprule
        \textbf{Παράμετρος} & \textbf{Τιμή} \\
        \midrule
        Συμμετέχοντες          & 30 \\
        Ηλικιακό εύρος         & 22–79 ετών \\
        Περιβάλλον             & Πραγματικές συνθήκες (AAL) \\
        Αισθητήρες             & Επιταχυνσιόμετρο + Γυροσκόπιο \\
        Συχνότητα δειγματοληψίας & 50\,Hz \\
        Αριθμός δειγμάτων (train) & 4252 \\
        Αριθμός δειγμάτων (test)  & 1492 \\
        Σύνολο δειγμάτων       & 5744 \\
        Αριθμός κλάσεων        & 6 \\
        Αριθμός χαρακτηριστικών & 561 \\
        \bottomrule
    \end{tabular}
\end{table}

Οι έξι κλάσεις δραστηριότητας είναι ίδιες με το UCI HAR (Walking,
Walking Upstairs, Walking Downstairs, Sitting, Standing, Laying),
καθιστώντας άμεση τη σύγκριση μεταξύ των δύο datasets. Οι βασικές
διαφορές που καθιστούν το AAL dataset ουσιαστικά πιο απαιτητικό
συνοψίζονται στον Πίνακα~\ref{tab:datasets_comparison}.

\begin{table}[H]
    \centering
    \caption{Σύγκριση UCI HAR και AAL Dataset.}
    \label{tab:datasets_comparison}
    \begin{tabular}{lll}
        \toprule
        \textbf{Χαρακτηριστικό} & \textbf{UCI HAR} & \textbf{AAL} \\
        \midrule
        Περιβάλλον συλλογής & Εργαστήριο & Πραγματικό \\
        Ηλικιακό εύρος      & 19–48 ετών & 22–79 ετών \\
        Θόρυβος σήματος     & Χαμηλός & Υψηλός \\
        Non-IID ετερογένεια & Μέτρια & Υψηλή \\
        Ρόλος στην εργασία  & Baseline & Κύριο πείραμα \\
        \bottomrule
    \end{tabular}
\end{table}

Η ευρύτερη ηλικιακή ποικιλομορφία των 22–79 ετών ενισχύει την
ετερογένεια των κινηματικών μοτίβων μεταξύ subjects, εντείνοντας
το Non-IID πρόβλημα και καθιστώντας το AAL dataset πλησιέστερο
σε πραγματικά σενάρια AAL παρακολούθησης.
